\section{Motivation}

% -----------
% FIRST FRAME

\begin{frame}
  \frametitle{Motivation}
  
  Autosubst is one of many libraries
  \footnote{We will be focusing on the Autosubst 1 library, but many of the limitations mentioned are common across these tools.}
  designed to facilitate much of the bureocratic work behind the formalization of calculi with binders, automating:
  \begin{itemize}
  \item<2-> Definitions of capture-avoinding substitution operations
  \item<3-> Proofs using substituion lemmas  
  \end{itemize}

  \onslide<4->{ However, such automation features only work when the calculus and substitution operations are fairly simple. }

\end{frame}

% ------------
% SECOND FRAME

\begin{frame} % [fragile]
  \frametitle{Various syntactic classes}
  
  One of the immediate limitations of Autosubst 1 is the use of calculi with many syntactic classes (which Autosubst 2 attemps to resolve).

  \onslide<2->{
    \begin{example}[System $\olb$]
      \vspace{-1.5em} 
      \begin{align*}
        \text{(terms)} \hspace{-5em} &&\quad t,u &::= \lb x.t \mid
                                                   \only<1-3>{ \der{x}{l} }
                                                   \only<4->{ {\color{red} \der{x}{l}} }  \mid tl \\
        \text{(lists)} \hspace{-5em} &&\quad l   &::= \nil \mid \cons{u}{l}    
      \end{align*}
      \vspace{-1.5em}
    \end{example}}
  
  \onslide<3->{ In our case, a simple workaround is to use polymorphic list types. }
  
  \onslide<4->{ \vspace{-1em} \seprule \vspace{-1em} }
  
  \onslide<4->{ Another unusual aspect of the system above is that the variable does not exists \emph{per se}. }

  \onslide<5->{ This forces a cut ($tl$) constructor on the substitution equation: \[ [u/x](\der{x}{l}) = u([u/x]l). \] }
  
\end{frame}

% ------------
% THIRD FRAME

\begin{frame}
  \frametitle{Various syntactic classes (cont.)}
  
  \begin{example}[System $\ulb$]
    \vspace{-1em} 
    \begin{align*}
      \text{(terms)} \hspace{-5em} &&\quad P,Q      &::= \lb x.t \mid \app{E} \\
      \text{(eliminations)} \hspace{-5em} &&\quad E &::= \only<1>{ x } \only<2->{ {\color{red} x} } \mid \hd{P} \mid EQ
    \end{align*}
    \vspace{-1em}
  \end{example}
  
  \onslide<2->{ Here, the variable lives in the elimination class and is not a term. }

  \onslide<3->{ Therefore, a similar forcing of constructor occurs on the substitution equation: \[ [P/x]x = {\color{red} hd}(P). \] }
  
\end{frame}

% ------------
% FOURTH FRAME

\begin{frame}
  \frametitle{Hereditary substitutions}

  Another kind of unsupported substitution operation is the hereditary substitution.

  Take as an example the $\beta$-normal terms of the ordinary $\lb$-calculus in vectorial notation,
  \[ M, N ::= \lb x.M \mid x \vec{M},\]
  and a closed substitution that normalizes by evaluation
  \[ \sub{N}{x}{(x\vec{M})} = \sub{(\lb y.N)}{x}{(xM\vec{M})} . \]
  
\end{frame}

%%% Local Variables:
%%% mode: LaTeX
%%% TeX-master: "../presentation"
%%% End:
